% year: תשפ"ה
% type: מבחן
% semester: א
% moed: מיוחד
% number: 2
% points: 21
% categories: func-analysis, func-limits
% source: exams/מבחנים/תשפה/מועד ג תשפה.pdf

\begin{question}
נגדיר פונקציה
\[ f(x) = \frac{x}{x^4 + 3} \]
\begin{enumerate}
  \item (11 נק') חשבו נקודות קיצון מקומי, ותחומי עלייה וירידה של הפונקציה בתחום הגדרתה. נמקו.
  \item (10 נק') חשבו את כל האסימפטוטות של הפונקציה. נמקו.
\end{enumerate}
\end{question}

\begin{solution}
תחום ההגדרה: $x^4 + 3 \ge 3 > 0$ לכל $x$, ולכן $f$ מוגדרת, רציפה וגזירה על כל $\R$.
\begin{enumerate}
  \item לפי כלל המנה:
  \[ f'(x) = \frac{(x^4 + 3) - x\cdot 4x^3}{(x^4+3)^2} = \frac{3 - 3x^4}{(x^4 + 3)^2} = \frac{3(1 - x^2)(1 + x^2)}{(x^4+3)^2}. \]
  המכנה ו-$1 + x^2$ חיוביים, ולכן סימן $f'$ הוא סימן $1 - x^2$:
  \begin{itemize}
    \item $|x| < 1$: $f' > 0$, $f$ עולה ב-$[-1, 1]$;
    \item $|x| > 1$: $f' < 0$, $f$ יורדת ב-$(-\infty, -1]$ וב-$[1, \infty)$.
  \end{itemize}
  לפי מבחן הנגזרת הראשונה: $x=-1$ נקודת מינימום מקומי, $f(-1) = -\frac14$; $x = 1$ נקודת מקסימום מקומי, $f(1) = \frac14$.
  (מכיוון ש-$f\to0$ באינסוף, אלה גם הקיצונים המוחלטים.)

  \item \textbf{אנכיות:} $f$ רציפה על כל $\R$, ולכן בכל נקודה $x_0$ מתקיים $\lim_{x\to x_0}f(x) = f(x_0)$ סופי, ואין אסימפטוטות אנכיות.

  \textbf{אופקיות/משופעות:}
  \[ \lim_{x\to\pm\infty}\frac{x}{x^4 + 3} = \lim_{x\to\pm\infty}\frac{1/x^3}{1 + 3/x^4} = 0. \]
  לכן $y = 0$ אסימפטוטה אופקית גם ב-$+\infty$ וגם ב-$-\infty$ (ולכן אין אסימפטוטה משופעת נוספת: $m = \lim \frac{f(x)}{x} = 0$, $n = \lim f(x) = 0$).

  \textbf{תשובה:} האסימפטוטה היחידה היא $y = 0$ (אופקית, בשני הכיוונים); אין אסימפטוטות אנכיות.
\end{enumerate}
\end{solution}
